\documentclass[10pt,a4paper]{amsart}
\usepackage[margin=19mm]{geometry}
\usepackage{amsmath,amssymb,mathtools}
\usepackage[scr=rsfs,cal=euler]{mathalfa}
\usepackage{xfrac}
\usepackage[unicode,hidelinks,bookmarks=false]{hyperref}
\newcommand{\ie}{i.e.\ }
\newcommand{\cf}{cf.\ }
\newcommand{\resp}{respectively\ }
\newcommand{\Subsection}{\subsection}
\providecommand{\nobreakdash}{\nobreak\hbox{-}}
\setlength{\emergencystretch}{2em}
\multlinegap0pt
\usepackage{amssymb}
\usepackage[shortlabels]{enumitem}
\usepackage[normalem]{ulem}
\newtheorem{lemma}{Lemma}
\title{High-Resolution Inertial Dynamics with Time-Rescaled Gradients for Nonsmooth Convex Optimization}
\author{Manh Hung Le, Andrea Simonetto}
\date{Source: arXiv:2603.25401v1 (CC BY 4.0)}
\begin{document}
\maketitle

\noindent\emph{Source and licence.} High-Resolution Inertial Dynamics with Time-Rescaled Gradients for Nonsmooth Convex Optimization. Version arXiv:2603.25401v1 (CC BY 4.0), distributed by arXiv under the Creative Commons Attribution 4.0 International licence, http://creativecommons.org/licenses/by/4.0/, as recorded in the arXiv licence metadata for this exact version and shown on the abstract page https://arxiv.org/abs/2603.25401v1. Full licence notice: \texttt{licenses/arxiv-2603.25401-CC-BY-4.0.txt}.\par\smallskip

\noindent\textit{Editorial scope.} Counted unit: the article's lemma moreau (source0.tex lines 1425--1435). The article's complete original proof of it is delivered with it (source0.tex lines 1437--1537, one \texttt{proof} environment of 101 lines). No mathematics is written, repaired or re-derived: the statement and the proof are the article's own lines, in the article's own notation, and no retained line is substituted or re-flowed.  No further source-local context is needed: every cross-reference of every retained line resolves inside the two delivered spans.  Nothing was omitted from any retained span, so the package carries no omission record.  No retained line contains a citation, so the wrapper carries no bibliography and no bibliography entry.  No retained line contains a figure environment, an image-inclusion command or drawing code, so no image is delivered and the package declares no asset dependency.  Editorial formatting only, in addition: an AMS \texttt{amsart} preamble that restates the article's own declarations and macros used by the retained lines, namely the environments \texttt{lemma} on separate counters, together with the standard packages those lines need.  The retained environments print in this document as Lemma 1 in order of appearance, whereas the article prints the same items with the numbering of its own full document.  The complete native source of the article as distributed in the arXiv e-print for this exact version is bundled unchanged as \texttt{sources/197207/source0.tex}.
\begin{lemma}\label{lem:derivative of moreau}
Let $\mathcal H$ be a real Hilbert space and let $A:\mathcal H\rightrightarrows\mathcal H$ be maximally monotone.
For $\gamma>0$ denote the resolvent $J_{\gamma A}=(I+\gamma A)^{-1}$ and the Yosida approximation $A_\gamma=\frac{1}{\gamma}(I-J_{\gamma A})$.
Let $\delta,\gamma:[t_0,+\infty)\to(0,+\infty)$ and $x:[t_0,+\infty)\to\mathcal H$ be $C^1$. Then, for any $x^{\star}\in\mathrm{zer}\,A$ and all $t\ge t_0$,
\[
\Big\Vert \tfrac{\mathrm{d}}{\mathrm{d}t}\big[\delta(t)A_{\gamma(t)}(x(t))\big] \Big\Vert
\le \frac{\delta(t)}{\gamma(t)}\Vert \dot{x}(t) \Vert
+ \Big[\frac{2\delta(t)\vert \dot{\gamma}(t)\vert}{\gamma(t)^2}
+ \frac{\vert \dot{\delta}(t)\vert}{\gamma(t)}\Big]\Vert x(t)-x^* \Vert .
\]
\end{lemma}

\begin{proof}
We use standard facts: for each $\gamma>0$, $J_{\gamma A}$ is firmly nonexpansive, and $A_\gamma$ is $1/\gamma$-Lipschitz and $\gamma$-cocoercive, i.e.,
\begin{equation}\label{eq:Lip-coco}
\|A_\gamma x-A_\gamma y\|\le \frac{1}{\gamma}\|x-y\|, \qquad
\langle A_\gamma x-A_\gamma y,\,x-y\rangle \ge \gamma\|A_\gamma x-A_\gamma y\|^2 .
\end{equation}

\textbf{Step 1: A basic bound via a zero of $A$.}
Fix $x^{\star}\in\mathrm{zer}\,A$ (so $J_{\gamma A}x^{\star}=x^{\star}$ and $A_\gamma(x^{\star})=0$). Apply \eqref{eq:Lip-coco} with $y=x^{\star}$ and use Cauchy--Schwarz inequality:
\begin{equation}
\gamma\|A_\gamma(x)\|^2 \le \langle A_\gamma(x),x-x^{\star}\rangle
\le \|A_\gamma(x)\|\,\|x-x^{\star}\| \;\Rightarrow\;
{\;\|A_\gamma(x)\|\le \frac{1}{\gamma}\|x-x^{\star}\|.\;}
\label{eq:basic-bound}
\end{equation}

\textbf{Step 2: Resolvent identity and parameter perturbation.}
We first show the resolvent identity: for any $\alpha,\beta>0$ and $x\in\mathcal H$,
\begin{equation}\label{eq:RI}
{\quad
J_{\alpha A}
= J_{\beta A}\!\left(\tfrac{\beta}{\alpha}I+\Big(1-\tfrac{\beta}{\alpha}\Big)J_{\alpha A}\right).
\quad}
\end{equation}
Indeed, let $u=J_{\alpha A}x$, so $(x-u)/\alpha\in Au$. Multiplying by $\beta$ gives
$(\beta/\alpha)(x-u)\in \beta Au$, hence
\[
u=(I+\beta A)^{-1}\!\Big(\tfrac{\beta}{\alpha}x+\Big(1-\tfrac{\beta}{\alpha}\Big)u\Big)
=J_{\beta A}\!\Big(\tfrac{\beta}{\alpha}x+\Big(1-\tfrac{\beta}{\alpha}\Big)u\Big),
\]
which is \eqref{eq:RI}.

From \eqref{eq:RI} we derive the parameter perturbation bound: for $\alpha,\beta>0$ and all $x$,
\begin{equation}\label{eq:PP}
{\quad
\|A_\beta(x)-A_\alpha(x)\|
\le \frac{2|\beta-\alpha|}{\alpha}\,\|A_\beta(x)\| .
\quad}
\end{equation}
To see this, set $u=J_{\alpha A}x$ and $v=J_{\beta A}x$. Using \eqref{eq:RI} with $\alpha,\beta$ swapped gives
\[
v=J_{\alpha A}\!\Big(\tfrac{\alpha}{\beta}x+\Big(1-\tfrac{\alpha}{\beta}\Big)v\Big).
\]
By nonexpansiveness of $J_{\alpha A}$,
\begin{equation}\label{eq:v-u}
\|v-u\|
\le \left\|\Big(\tfrac{\alpha}{\beta}-1\Big)x+\Big(1-\tfrac{\alpha}{\beta}\Big)v\right\|
= \frac{|\alpha-\beta|}{\beta}\,\|x-v\|.
\end{equation}
Now
\[
A_\beta(x)-A_\alpha(x)
= \frac{x-v}{\beta}-\frac{x-u}{\alpha}
= \Big(\frac{1}{\beta}-\frac{1}{\alpha}\Big)(x-v) + \frac{1}{\alpha}(v-u),
\]
whence, using \eqref{eq:v-u},
\[
\|A_\beta(x)-A_\alpha(x)\|
\le \frac{|\alpha-\beta|}{\alpha\beta}\|x-v\|
+ \frac{1}{\alpha}\cdot \frac{|\alpha-\beta|}{\beta}\|x-v\|
= \frac{2|\alpha-\beta|}{\alpha\beta}\|x-v\|.
\]
Since $\|x-v\|/\beta=\|A_\beta(x)\|$, we obtain \eqref{eq:PP}.

\textbf{Step 3: Difference quotient and limit.}
Define $g(t):=\delta(t)A_{\gamma(t)}(x(t))$. For small $h\neq0$,
\[
\begin{aligned}
\|g(t+h)-g(t)\|
&= \|\delta(t+h)A_{\gamma(t+h)}(x(t+h))-\delta(t)A_{\gamma(t)}(x(t))\|\\
&\le \underbrace{\delta(t)\|A_{\gamma(t)}(x(t+h))-A_{\gamma(t)}(x(t))\|}_{\mathrm{(I)}}\\
&\quad + \underbrace{\delta(t)\|A_{\gamma(t+h)}(x(t+h))-A_{\gamma(t)}(x(t+h))\|}_{\mathrm{(II)}}\\
&\quad + \underbrace{|\delta(t+h)-\delta(t)|\,\|A_{\gamma(t+h)}(x(t+h))\|}_{\mathrm{(III)}}.
\end{aligned}
\]
Using the $1/\gamma$-Lipschitzness in $x$,
\[
\mathrm{(I)} \le \frac{\delta(t)}{\gamma(t)}\,\|x(t+h)-x(t)\|.
\]
Using \eqref{eq:PP} with $\alpha=\gamma(t)$ and $\beta=\gamma(t+h)$,
\[
\mathrm{(II)} \le \frac{2\delta(t)|\gamma(t+h)-\gamma(t)|}{\gamma(t)}\,\|A_{\gamma(t+h)}(x(t+h))\|.
\]
Applying \eqref{eq:basic-bound} at $(x,\gamma)=(x(t+h),\gamma(t+h))$ yields
\[
\mathrm{(II)} \le \frac{2\delta(t)|\gamma(t+h)-\gamma(t)|}{\gamma(t)\gamma(t+h)}\,\|x(t+h)-x^*\|.
\]
Similarly, by \eqref{eq:basic-bound},
\[
\mathrm{(III)} \le \frac{|\delta(t+h)-\delta(t)|}{\gamma(t+h)}\,\|x(t+h)-x^*\|.
\]

Divide by $|h|$ and let $h\to0$. Since $x,\gamma,\delta$ are $C^1$ and $(x,\gamma)\mapsto A_\gamma(x)$ is continuous, we obtain
\[
\Big\Vert \tfrac{\mathrm{d}}{\mathrm{d}t}\big[\delta(t)A_{\gamma(t)}(x(t))\big] \Big\Vert
\le \frac{\delta(t)}{\gamma(t)}\Vert \dot{x}(t) \Vert
+ \Big[\frac{2\delta(t)\vert \dot{\gamma}(t)\vert}{\gamma(t)^2}
+ \frac{\vert \dot{\delta}(t)\vert}{\gamma(t)}\Big]\Vert x(t)-x^* \Vert,
\]
as claimed.
\end{proof}

\end{document}
